\chapter{Peeling Skin}

\section*{Definition}
Desquamation or shedding of the outer epidermal layer of skin, which may be localized or generalized, resulting from environmental damage, infection, inflammation, or systemic disease.

\section*{Pathophysiology}
Peeling skin results from disruption of intercellular adhesion between keratinocytes or accelerated epidermal turnover. Sunburn causes direct keratinocyte DNA damage triggering apoptosis and detachment. In inflammatory conditions (eczema, psoriasis), dysregulated immune responses accelerate epidermal proliferation and shedding. Infectious causes (staphylococcal scalded skin syndrome, fungal infections) directly disrupt desmosomal adhesion.

\section*{Causes}
\begin{itemize}
  \item Sunburn (UV-induced skin damage)
  \item Dry skin (xerosis) from low humidity or over-washing
  \item Allergic contact dermatitis (poison ivy, nickel, cosmetics)
  \item Atopic dermatitis (eczema)
  \item Psoriasis (plaque or guttate)
  \item Fungal infections (tinea pedis/athlete's foot)
  \item Kawasaki disease (in children, with fever and rash)
  \item Scarlet fever (desquamation after fever resolves)
  \item Staphylococcal scalded skin syndrome (emergency)
  \item Medication reactions (Stevens-Johnson syndrome — emergency)
\end{itemize}

\section*{When to See a Doctor}
See your doctor if:
\begin{itemize}
  \item Severe or worsening symptoms
  \item Peeling skin with fever, blistering, mucosal involvement, or widespread skin detachment (possible SJS/TEN)
  \item Accompanied by other concerning signs
\end{itemize}

\section*{Related Symptoms}
\begin{itemize}
  \item Rash
  \item Itching
  \item Dry skin
  \item Redness
  \item Blisters
\end{itemize}

\textbf{Important:} If this symptom persists, your doctor will check for serious underlying causes.

\section*{Self-Care / Home Management}
\begin{itemize}
  \item Rest and avoid activities that make it worse.
  \item Maintain adequate hydration and a balanced diet.
  \item Over-the-counter remedies may provide symptomatic relief (consult a pharmacist or doctor).
  \item Monitor symptoms and seek professional advice if they persist or worsen.
  \item Avoid self-medicating with prescription medications without medical guidance.
\end{itemize}

\section*{How Your Doctor Will Evaluate You}
\begin{itemize}
  \item A thorough history and physical examination are the first steps in evaluation.
  \item Laboratory tests (blood work, urinalysis) may be ordered based on clinical suspicion.
  \item Imaging studies (X-ray, ultrasound, CT, MRI) may be indicated when structural pathology is suspected.
  \item Specialist referral is arranged when the diagnosis remains uncertain or requires specific expertise.
\end{itemize}

\section*{Treatment Options}
\begin{itemize}
  \item Treatment targets the underlying cause identified through diagnostic evaluation.
  \item Supportive care and symptom management are provided while awaiting definitive therapy.
  \item Medications, lifestyle modifications, and physical therapies may be prescribed as appropriate.
  \item Follow-up ensures treatment effectiveness and allows timely adjustments to the care plan.
\end{itemize}

\clearpage